\documentclass[11pt, a4paper, oneside]{Thesis} % Paper size, default font size and one-sided paper
\usepackage{floatrow}
\floatsetup[table]{capposition=top}
\usepackage{textcomp}
\usepackage{wrapfig}
\usepackage{lscape}
\usepackage{rotating}
\usepackage{hyperref}
\usepackage{graphicx}
\usepackage{algorithm}
\usepackage{algorithmic}
%\usepackage{algpseudocode}
\usepackage{float}
\usepackage{caption}
\usepackage{amsmath}
\usepackage{upgreek}
\usepackage{gensymb}
\usepackage{csquotes}
\usepackage{pdfpages}
\usepackage{lipsum}  
%\usepackage[open]{bookmark}

% acronyms
\usepackage{acronym}
\usepackage{array}
\newcolumntype{L}{>{\centering\arraybackslash}m{3cm}}

%\usepackage[acronym]{glossaries}
% prints author names as small caps


\makeatletter
\AtBeginDocument{%
  \renewcommand*{\AC@hyperlink}[2]{%
    \begingroup
      \hypersetup{hidelinks}%
      \hyperlink{#1}{#2}%
    \endgroup
  }%
}
\makeatother

%\usepackage{subcaption} %incompatible with subfig
\graphicspath{{Pictures/}} % Specifies the directory where pictures are stored
\usepackage[square, numbers]{natbib} % Use the natbib reference package - read up on this to edit the reference style; if you want text (e.g. Smith et al., 2012) for the in-text references (instead of numbers), remove 'numbers' v

\hypersetup{urlcolor=black, colorlinks=true} % Colors hyperlinks in blue - change to black if annoyingv`
\title{\ttitle} % Defines the thesis title - don't touch this

\begin{document}
\makeatletter
\renewcommand*{\NAT@nmfmt}[1]{\textsc{#1}}
\makeatother

% prints author names as small caps


\frontmatter % Use roman page numbering style (i, ii, iii, iv...) for the pre-content pages

\setstretch{1.6} % Line spacing of 1.6 (double line spacing)

% Define the page headers using the FancyHdr package and set up for one-sided printing
\fancyhead{} % Clears all page headers and footers
\rhead{\thepage} % Sets the right side header to show the page number
\lhead{} % Clears the left side page header

\pagestyle{fancy} % Finally, use the "fancy" page style to implement the FancyHdr headers

\newcommand{\HRule}{\rule{\linewidth}{0.5mm}} % New command to make the lines in the title page

% PDF meta-data
\hypersetup{pdftitle={\ttitle}}
\hypersetup{pdfsubject=\subjectname}
\hypersetup{pdfauthor=\authornames}
\hypersetup{pdfkeywords=\keywordnames}

%----------------------------------------------------------------------------------------
%	TITLE PAGE
%----------------------------------------------------------------------------------------

\begin{titlepage}
\begin{center}

\HRule \\[0.4cm] % Horizontal line
{\huge \bfseries \ttitle}\\[0.4cm] % Thesis title
\HRule \\[1.5cm] % Horizontal line
 
\large \textit{A deep dive in to the QR algorithm using Golang}\\[0.3cm] % University requirement text
\textit{by}\\[0.4cm]
\authornames \\



\vfill
\graphicspath{ {./Figures/} }
\begin{figure}[hb]
  \centering
  \includegraphics[width=0.35\linewidth]{Pictures/IITHLOGO.png}
\end{figure}

\DEPTNAME\\ % Research group name and department name
\textsc{ \UNIVNAME}\\[1.5cm] % University name
\large \today\\[2cm] % Date


\end{center}

\end{titlepage}

%----------------------------------------------------------------------------------------
%	DECLARATION PAGE
%	Your institution may give you a different text to place here
%----------------------------------------------------------------------------------------

\Declaration{\addtocontents{toc}{\vspace{1em}}} % Add a gap in the Contents, for aesthetics
\setcounter{page}{2}

\begin{minipage}{\textwidth}
    
    It is certified that the work contained in this thesis entitled \textbf{\enquote{\ttitle}} by \textbf{\authornames} has been carried out under my supervision and that it has not been submitted elsewhere for a degree.
        
\end{minipage}

\vspace{20.00mm}

\begin{minipage}{0.49\textwidth}
		\begin{flushleft}
			{\supnameA\\ 
			Professor\\
			\deptname\\
			\univname}
		\end{flushleft}
\end{minipage}
\vfill
\clearpage % Start a new page
\StudentDeclaration{\addtocontents{toc}{\vspace{1em}}} % Add a gap in the Contents, for aesthetics

This is to certify that the thesis titled \textbf{``\ttitle''} has been authored by me. It presents the research conducted by me under the supervision of \textbf{\supnameA} .\par

To the best of my knowledge, it is an original work, both in terms of research content and narrative, and has not been submitted elsewhere, in part or in full, for a degree. Further, due credit has been attributed to the relevant state-of-the-art and collaborations with appropriate citations and acknowledgments, in line with established norms and practices.\\ [2cm]
\begin{minipage}{.5\textwidth}
		\begin{flushleft}
			{\authornames\\ Roll No. EE24BTECH11001 \\
			\normalsize{\href{http://www.iitk.ac.in/mse}{EE Department}\\
			\univname}}
		\end{flushleft}
\end{minipage}
\vfill

\clearpage % Start a new page
%----------------------------------------------------------------------------------------
%	ABSTRACT PAGE
%----------------------------------------------------------------------------------------

%\addtotoc{Abstract} % Add the "Abstract" page entry to the Contents
\lhead{\emph{Abstract}}

\abstract{\addtocontents{toc}{\vspace{1em}} % Add a gap in the Contents, for aesthetics
In general, any method for computing eigenvalues necessarily involves an infinite number of steps. This is because finding eigenvalues of an $n \times n$ matrix is equivalent to finding the roots of its characteristic polynomial of degree $n$ and Computing coefficients of characteristic polynomial requires computation of the determinant, however, the problem of finding the roots of a polynomial can be very ill-conditioned and for $n > 4$ such roots cannot be found (By Abel’s theorem), in general, in a finite number of steps. In other words, we must consider iterative methods producing, at step $k$, an approximated eigenvector $x_k$ associated with an approximated eigenvalue $\lambda_k$ that converge to the desired eigenvector x and eigenvalue $\lambda$ as the number of iterations becomes larger. Here I try to motivate the concept of the QR algorithm for finding eigenvalues of arbitrary matrices.

}

%----------------------------------------------------------------------------------------
%	Declaration Page
%----------------------------------------------------------------------------------------

%----------------------------------------------------------------------------------------
%	ACKNOWLEDGEMENTS
%----------------------------------------------------------------------------------------
\clearpage % Start a new page
\setstretch{1.3} % Reset the line-spacing to 1.3 for body text (if it has changed)

\clearpage % Start a new page

%----------------------------------------------------------------------------------------
%	LIST OF CONTENTS/FIGURES/TABLES PAGES
%----------------------------------------------------------------------------------------

\pagestyle{fancy} % The page style headers have been "empty" all this time, now use the "fancy" headers as defined before to bring them back

\lhead{\emph{Contents}} % Set the left side page header to "Contents"
\tableofcontents % Write out the Table of Contents

%----------------------------------------------------------------------------------------
%	ABBREVIATIONS
%----------------------------------------------------------------------------------------

\clearpage % Start a new page

\setstretch{1.5} % Set the line spacing to 1.5, this makes the following tables easier to read
%+\begin{acronym}[XXXXXXXXX] % Give the longest label here so that the list is nicely aligned
%\acro{EBSD}{electron backscatter diffraction}
% \acro{SEM}{scanning electron microscope}
% \acro{WHA}{tungsten heavy alloys}
% \acro{SPS}{spark plasma sintering}
% \acro{ND}{normal direction to the rolling plane}
% \acro{CD}{compression direction}
% \acro{RD}{rolling direction in a rolled plate}
% \acro{LD}{Loading direction}
% \acro{TD}{transverse direction in a rolled plate}
% \acro{XRD}{X-ray diffraction}
% \acro{ODF}{orientation distribution function}
% \acro{OM}{optical microscope}
% \acro{IQ}{image quality}
% \acro{IPF}{inverse pole figure}
% \acro{FESEM}{field emission scanning electron microscopy}
% \acro{f.c.c.}{face centred cubic}
% \acro{b.c.c.}{body centred cubic}
% \acro{EPMA}{electron probe microanalyser}
% \acro{EDS}{energy dispersive spectrometer}
% \acro{LPS}{liquid phase sintering}
% \acro{FFT}{fast Fourier transform}
% \acro{HEM}{homogeneous effective medium}
% \acro{RVE}{representative volume element}
% \acro{CRSS}{critical resolved shear stress}
% \acro{PF}{pole figure}
% \acro{VPSC}{viscoplastic self-consistent}
% \acro{STECA}{Stress Texture Calculator}
% \acro{DAMASK}{D\"{u}sseldorf Advanced Material Simulation Kit}
% \acro{KAM}{kernal average misorientation}
% \acro{GND}{geometry necessary dislocation}
% \acro{GROD}{grain reference orientation deviation}
% \acro{CAD}{computer-aided design}
% \acro{AVACS}{average number of active systems per grain}
% \acro{SRO}{short range ordering}
% \acro{TWIP}{twinning induced plasticity}
% \acro{TRIP}{transformation induced plasticity}
% \acro{h.c.p.}{hexagonal close packed}
% \acro{YS}{yield strength}
% \acro{UTS}{ultimate tensile strength}
% \acro{MSBs}{micro-shear bands}
% \acro{DB}{deformation bands}
% \acro{TB}{transition bands}
% \acro{SFE}{stacking fault energy}
% \acro{SHPB}{split-Hopkinson pressure bar}
% \acro{ppm}{parts per million}
% \acro{CES}{Cambridge Engineering Selector}
% \acro{DBTT}{ductile-to-brittle transition temperature}
% \acro{mrd}{multiples of a random distribution}
% \acro{FWHM}{full width at half maximum}
% \acro{CPFEM}{crystal plasticity finite element methods}

\end{acronym}

%----------------------------------------------------------------------------------------
%	PHYSICAL CONSTANTS/OTHER DEFINITIONS
%----------------------------------------------------------------------------------------
%
% \clearpage % Start a new page

% \lhead{\emph{Physical Constants}} % Set the left side page header to "Physical Constants"

% \listofconstants{lrcl} % Include a list of Physical Constants (a four column table)
% {
% Speed of Light & $c$ & $=$ & $2.997\ 924\ 58\times10^{8}\ \mbox{ms}^{-\mbox{s}}$ (exact)\\
% % Constant Name & Symbol & = & Constant Value (with units) \\
% }

%----------------------------------------------------------------------------------------
%	SYMBOLS
%----------------------------------------------------------------------------------------

\clearpage % Start a new page
% \\[2cm]


%\vfill{}


\clearpage % Start a new pag


\setstretch{1.3} % Return the line spacing back to 1.3
%
\pagestyle{empty} % Page style needs to be empty for this page
%

%\addtocontents{toc}{\vspace{2em}} % Add a gap in the Contents, for aesthetics

%----------------------------------------------------------------------------------------
%	THESIS CONTENT - CHAPTERS
%----------------------------------------------------------------------------------------

\mainmatter % Begin numeric (1,2,3...) page numbering

\pagestyle{fancy} % Return the page headers back to the "fancy" style

% Include the chapters of the thesis as separate files from the Chapters folder
% Uncomment the lines as you write the chapters

% Chapter Template

\chapter{Introduction} % Main chapter title

\label{Chapter1}

\lhead{Chapter 1. \emph{Introduction}} % Change X to a consecutive number; this is for the header on each page - perhaps a shortened title

%----------------------------------------------------------------------------------------
%	SECTION 1
%----------------------------------------------------------------------------------------

\section{Why is this hard?}
    Symmetry makes life less complex. This statement is especially true when applied to the problem at hand. The problem of finding eigenvalues is already quite involved in the case of real and symmetric matrices. At least, these symmetric matrices have a guarantee of all eigenvalues being real. In the general case of complex matrices, it becomes mandatory for us to use iterative methods to "converge" to the true eigenvalues.
    \\
    The spectrum of algorithms available on the mathematical market can be classified broadly into two categories :
    \begin{enumerate}
        \item \textbf{Direct Methods:} Direct methods are typically used on dense matrices and cost
$O(n^3)$ operations to compute all eigenvalues and eigenvectors; this cost is relatively insensitive to the actual matrix entries.
        \item \textbf{Iterative Methods: }
         These are usually applied
to sparse matrices or matrices for which matrix-vector multiplication is the
only convenient operation to perform. Iterative methods typically provide
approximations only to a subset of the eigenvalues and eigenvectors and are
usually run only long enough to get a few adequately accurate eigenvalues
rather than a large number. Their convergence properties depend strongly on
the matrix entries.
    \end{enumerate}
    Here I will try and motivate one of the most often used direct algorithms for the task at hand - namely, the QR algorithm.

\section{Design Considerations}
Before we set out to accomplish our goal to calculate eigenvalues for arbitrary matrices, we must put some thought behind what qualities our "ideal" algorithm should have :
\begin{enumerate}
    \item \textbf{Accuracy:} We wouldn't want to waste time and effort implementing an algorithm which doesn't provide the true value of eigenvalues up to a desired precision.

    \item \textbf{Efficiency:} While we need to come to terms with the fact that we can't reap the efficiency rewards of exploiting structure of specific types of matrices, we should strive for the fastest algorithm which gets the job done.
    \item \textbf{Convergence:} One might argue that this is a subset of the previous design consideration, but it is important enough to point out for itself. Oceans worth of ink has been used to improve convergence properties of various algorithms, so we must also look into this aspect.
    \item \textbf{Stability:} A slightly more nuanced design consideration. But we want the algorithm not to blow up if we change the entries within the matrix by a relatively small amount.

\end{enumerate}

% Chapter Template

\chapter{The Basic QR Algorithm}

\label{Chapter2}

\lhead{Chapter 2. \emph{The Basic QR Algorithm}} 

\section{Preliminaries}
\label{Preliminaries}
\subsection{Similarity Transformations}
    \textbf{Definition: } Two matrices $A, B \in \mathbb{C}^{n \times n}$ are said to be similar if there exists an invertible matrix $P \in \mathbb{C}^{n \times n}$ such that
    \begin{align}
        B = P^{-1}AP
    \end{align}
    Similar matrices represent the same linear map under two (possibly) different bases, with P being the change-of-basis matrix. Because matrices are similar if and only if they represent the same linear operator with respect to (possibly) different bases, similar matrices share all properties of their shared underlying operator like rank, 
    the characteristic polynomial and by association the eigenvalues and the eigenvectors.
\subsection{Unitary Matrices}
    An invertible complex square matrix $U$ is unitary if its matrix inverse $U^{-1}$ equals its conjugate transpose $U^*$, that is, if
    \begin{align}
        UU^* = U^*U = I
    \end{align}
    where $I$ is the identity matrix.
\subsection{Schur Decomposition}
    It allows one to write an arbitrary complex square matrix as unitarily similar to an upper triangular matrix whose diagonal elements are the eigenvalues of the original matrix. If $A \in \mathbb{C}^{n \times n}$ is an $n \times n$ square matrix with complex entries, then A can be expressed as
    \begin{align}
        A = QUQ^{-1}
    \end{align}
    for some unitary matrix $Q$ (so that the inverse $Q^{-1}$ is also the conjugate transpose $Q^*$ of $Q$), and some upper triangular matrix $U$. Since $U$ is similar to $A$, it has the same spectrum, and since it is triangular, \textbf{its eigenvalues are the diagonal entries of $U$}. It must be noted that all square have a Schur decomposition, although it may not be unique.
\subsection{Hessenberg Form}
A Hessenberg matrix is a special kind of square matrix, one that is "almost" triangular. To be exact, an upper Hessenberg matrix has zero entries below the first sub-diagonal, and a lower Hessenberg matrix has zero entries above the first super-diagonal.
\begin{align}
    \begin{bmatrix}
        \times & \times & \times& \times& \times\\
        \times & \times & \times& \times& \times\\
        0 & \times & \times& \times& \times\\
        0 & 0 & \times& \times& \times\\
        0 & 0 & 0& \times& \times\\
    \end{bmatrix}
\end{align}
\section{The Basic Idea Behind QR}
    Our primary goal is to get the Schur decomposition of the given matrix and read off the diagonal values to be the eigenvalues.\\

    The QR algorithm works off of the following observation,\\
    For every matrix $A$ we can represent it as the product of a unitary matrix $Q$ and an upper triangular matrix $R$.

    \begin{align}
        A_1 &= Q_1R_1\\
        A_2 &= R_1Q_1 = Q_1^*AQ_1 = Q_2R_2\\
        A_3 &= R_2Q_2 = Q_2^*Q^*1AQ_1Q_2\\
        \vdots\\
        A_{k+1} &= Q_{k-1}^*Q_{k-2}^*\cdots Q_2^*Q_1^*AQ_1Q_2\cdots Q_{k-2}Q_{k-1}
    \end{align}

    We can see that $A_{k+1}$ is similar to $A_k$. It can be shown that under certain conditions, as we run these iterations, $A_{k+1}$ will converge to an upper triangular matrix\footnote{We will explore this notion in Chapter 4}. So now if we take
    \begin{align}
        U = Q_1Q_2\cdots Q_{k-2}Q_{k-1}
    \end{align}
    then, 
    \begin{align}
        A = U^*A_{\infty}U
    \end{align}
    So we have successfully calculated the Schur decomposition of the desired matrix.

\section{Improving the QR Algorithm}
Two of the major drawbacks of this unoptimized algorithm are as follows:
\begin{enumerate}
    \item The convergence of the algorithm is slow. In fact it can be arbitrarily slow if eigenvalues are very close to each other.
\item The algorithm is expensive. Each iteration step requires the computation of the QR
factorization of a full $n \times n$ matrix, i.e., each single iteration step has a complexity
$O(n^3)$. Even if we assume that the number of steps is proportional to $n$ we would
get an $O(n^4)$ complexity. The latter assumption is not even assured, see point 1 of
this discussion.
\end{enumerate}
Because of the number of arithmetic operations involved, the QR algorithm is practical
only when applied to a matrix in upper Hessenberg form (i.e., $a_{ij} = 0$ for $i > j + 1$). Hence
the method consists of two steps. In the first step, we find a similar matrix which is in upper
Hessenberg form. The second step is then to apply the QR algorithm to this new matrix.
One can show that if A is in upper Hessenberg form, applying the QR algorithm results in
a sequence of matrices that are also in this form. 


            




 
% Chapter Template

\chapter{The Hessenberg QR Algorithm} % Main chapter title

\label{Chapter3} % Change X to a consecutive number; for referencing this chapter elsewhere, use \ref{ChapterX}

\lhead{Chapter 3. \emph{The Hessenberg QR Algorithm}} % Change X to a consecutive number; this is for the header on each page - perhaps a shortened title

The flow of this algorithm is as follows:
\begin{enumerate}
    \item Convert given matrix to upper Hessenberg form using Householder Reflectors in $O(n^3)$ time complexity.
    \item Use Givens Rotations to reduce the matrix to upper triangular matrix in $O(n^2)$ time.
\end{enumerate}
The motivation to convert to upper Hessenberg form is that the computations are much cheaper and hence faster, which achieves Design Consideration 2.
\section{Householder Reflectors}
\textbf{Definition: }A matrix of the form
\begin{align}
    P = I - 2\textbf{uu}^*\\
\end{align}
is called a Householder reflector.\\
It is easy to verify that Householder reflectors are Hermitian and that 
$P^2 = I$. From this
we deduce that P is unitary. It is clear that we only have to store the Householder
vector u to be able to multiply a vector (or a matrix) with P,
\begin{align}
    P\textbf{x} = \textbf{x} - \textbf{u}(2\textbf{u}^*\textbf{x})\\
\end{align}
A task that we repeatedly want to carry out with Householder reflectors is to transform
a vector \textbf{x} on a multiple of \textbf{$e_1$}
\begin{align}
    P\textbf{x} = \textbf{x} - \textbf{u}(2\textbf{u}^*\textbf{x}) = \alpha \textbf{$e_1$}\\
\end{align}
Since P is unitary, we must have $\alpha = \rho 
||\textbf{x}||$, where $\rho \in \mathbb{C}$ has absolute value one. Therefore,
\begin{align}
    \textbf{u} = \frac{\textbf{x} - \rho||\textbf{x}||\textbf{$e_1$}}{||\textbf{x} - \rho||\textbf{x}||\textbf{$e_1$}||}
\end{align}
We can freely choose $\rho$ provided that $|\rho| = 1$. Let $x_1 = |x_1|e^{i\textbf{$u_k$}^*\phi}$. To avoid numerical
cancellation we set $\rho = -e^{i\phi}$
.
In the real case, one commonly sets $\rho$ = -sign($x_1$). If $x_1 = 0$ we can set $\rho$ in any way.
\subsection{Algorithm to Convert to Hessenberg Form}
The main algorithm is as follows
\begin{algorithm}
\caption{An algorithm with caption}\label{alg:cap}
    \begin{algorithmic}
    \FOR {k = 1 to n - 2}
        \STATE Generate the Householder reflector $P_k$
        \STATE $A_{k+1:n, k:n} := A_{k+1:n, k:n} - \textbf{$u_k$}(2\textbf{$u_k$}^*A_{k+1:n, k:n})$
        \COMMENT{Apply $P_k$ to $A$ from left}
        \STATE $A_{1:n, k+1:n} := A_{1:n, k+1:n} - 2(A_{1:n, k+1:n}\textbf{$u_k$})\textbf{$u_k$}^*$
        \COMMENT{Apply $P_k$ to $A$ from right}
    \ENDFOR
   \end{algorithmic}
\end{algorithm}

To get a visualization of the transformation,
\begin{align}  
    P_1 = \begin{bmatrix}
        1 & 0 & 0 & 0 & 0 \\
        0 & \times & \times& \times& \times\\
        0 & \times & \times& \times& \times\\
        0 & \times & \times& \times& \times\\
        0 & \times & \times& \times& \times\\
    \end{bmatrix} = \begin{bmatrix}
        1 & \textbf{0}^{\top}\\
        \textbf{0} & I_4 - 2u_1u_1^*
    \end{bmatrix}    
\end{align}
\\
\begin{align}
    A = \begin{bmatrix}
        \times & \times & \times& \times& \times\\
        \times & \times & \times& \times& \times\\
        \times & \times & \times& \times& \times\\
        \times & \times & \times& \times& \times\\
        \times & \times & \times& \times& \times\\
    \end{bmatrix} \xrightarrow[P_1*]{} \begin{bmatrix}
        \times & \times & \times& \times& \times\\
        \times & \times & \times& \times& \times\\
        0 & \times & \times& \times& \times\\
        0 & \times & \times& \times& \times\\
        0 & \times & \times& \times& \times\\
    \end{bmatrix} \xrightarrow[*P_1]{}  \begin{bmatrix}
        \times & \times & \times& \times& \times\\
        \times & \times & \times& \times& \times\\
        0 & \times & \times& \times& \times\\
        0 & \times & \times& \times& \times\\
        0 & \times & \times& \times& \times\\
    \end{bmatrix}
\end{align}
The multiplication of $P_1$ from the left inserts the desired zeros in column 1 of $A$. The
multiplication from the right is necessary in order to have similarity. Because of the
nonzero structure of $P_1$ the first column of $P_1A$ is not affected. Hence, the zeros stay
there.
\\
The reduction continues in a similar way:
\begin{multline}
    P_1^*AP_1 = \begin{bmatrix}
        \times & \times & \times& \times& \times\\
        \times & \times & \times& \times& \times\\
        0 & \times & \times& \times& \times\\
        0 & \times & \times& \times& \times\\
        0 & \times & \times& \times& \times\\
    \end{bmatrix} \xrightarrow[P_2*/*P_2]{} \begin{bmatrix}
        \times & \times & \times& \times& \times\\
        \times & \times & \times& \times& \times\\
        0 & \times & \times& \times& \times\\
        0 & 0 & \times& \times& \times\\
        0 & 0 & \times& \times& \times\\
    \end{bmatrix} 
    \xrightarrow[P_3*/*P_3]{} \begin{bmatrix}
        \times & \times & \times& \times& \times\\
        \times & \times & \times& \times& \times\\
        0 & \times & \times& \times& \times\\
        0 & 0 & \times& \times& \times\\
        0 & 0 & 0& \times& \times\\
    \end{bmatrix}\\ = P_3P_2P_1AP_1P_2P_3
\end{multline}
\\
It is worth pointing out that each of the pairs of multiplications by $P_k$ on the right and left are similarity transformations, hence they have no effect on the eigenvalues of the underlying matrix.
\section{Givens Rotations}
Givens rotations are used to zero specific elements of a vector or matrix by rotating the vector in the plane of two coordinates. A Givens rotation matrix is defined as:
\[
G(i, j, \theta) =
\begin{bmatrix}
1 & \cdots & 0 & \cdots & 0 & \cdots & 0 \\
\vdots & \ddots & \vdots & \ddots & \vdots & \ddots & \vdots \\
0 & \cdots & c & \cdots & s & \cdots & 0 \\
\vdots & \ddots & \vdots & \ddots & \vdots & \ddots & \vdots \\
0 & \cdots & -\overline{s} & \cdots & \overline{c} & \cdots & 0 \\
\vdots & \ddots & \vdots & \ddots & \vdots & \ddots & \vdots \\
0 & \cdots & 0 & \cdots & 0 & \cdots & 1
\end{bmatrix},
\]
where 
\begin{align}
    c_k = \frac{\overline{x_i}}{\sqrt{|x_i|^2 + |x_j|^2}}\\
    s_k = \frac{\overline{x_j}}{\sqrt{|x_i|^2 + |x_j|^2}}   
\end{align}
Now the algorithm to run one iteration for the upper Hessenberg form is as follows:
\begin{algorithm}
\caption{A Hessenberg QR Step}
\begin{algorithmic}[1]
\FOR{\(k = 1, 2, \dots, n-1\)}
    \STATE Generate \(G_k\) and then apply it: \(H := G(k, k+1, \theta_k)^* H\).
    \STATE Compute: \([c_k, s_k] := \text{givens}(H_{k,k}, H_{k+1,k})\).
    \STATE Update rows \(k:k+1\): \(H_{k:k+1,k:n} = \begin{bmatrix} c_k & -s_k \\ s_k & c_k \end{bmatrix} H_{k:k+1,k:n}\).
\ENDFOR
\FOR{\(k = 1, 2, \dots, n-1\)}
    \STATE Apply the rotations \(G_k\) from the right:
    \STATE \(H_{1:k+1,k:k+1} = H_{1:k+1,k:k+1} \begin{bmatrix} c_k & s_k \\ -s_k & c_k \end{bmatrix}\).
\ENDFOR
\end{algorithmic}
\end{algorithm}

\section{Complexity}
We will now compute time complexities for both the major steps in our Hessenberg QR algorithm:
\begin{enumerate}
    \item \textbf{Computing Hessenberg :}
        \begin{enumerate}
            \item Application of $P_k$ from the left :
            $\sum_{k=1}^{n-2}4(n-k-1)(n-k) \approx \frac{4n^3}{3}$ 
            \item Application of $P_k$ from the right :
            $\sum_{k=1}^{n-2}4n(n-k) \approx 2n^3$
        \end{enumerate}

    \item \textbf{Running Givens Rotations: }
    Neglecting the minimal overhead of calculating $c_k$ and $s_k$, executing the 2 loops in the algorithm costs :
    \begin{align}
        2\times\sum_{i=1}^{n-1}6i = 6n(n-1)\approx 6n^2
    \end{align}
\end{enumerate}

As we can see, while running the iterations to converge to the Schur decomposition, the Hessenberg algorithm requires $6n^2$ arithmetic computations for each iteration while the naive QR algorithm takes $\frac{7n^3}{3}$ operations.


% Chapter Template

\chapter{Addressing Convergence Issues}

\label{Chapter4}

\lhead{Chapter 4. \emph{Addressing Convergence Issues}} 

\section{Does the Algorithm Always Converge?}
At the start of the discussion I talked about the convergence of the QR iteration to an upper triangular matrix "under certain conditions". We will now explore what these conditions are, and when does the convergence fail to occur. 
\subsection{Necessary Condition For Convergence}
The iterations will converge in a "reasonable" number of iterations if
\begin{enumerate}
    \item If all eigenvalues have distinct absolute values. Precisely, if $\lambda_1, \lambda_2 \cdots \lambda_n$ be the eigenvalues of the matrix, for convergence to occur
    \begin{align}
        |\lambda_1| < |\lambda_2|< \cdots < |\lambda_n|
    \end{align}
    \item If the ratio of any two eigenvalues is "reasonably"\footnote{See Reference No. 2 for more elaborate discussion} large. Precisely,
    \begin{align}
        \frac{\lambda_i}{\lambda_j} >> 1, i \ne j
    \end{align}
\end{enumerate}

\subsection{What Happens if Absolute Values are not Unique?}
The short answer to this question is that you will get $2 \times 2$ blocks called \textbf{Jordan Blocks} where there is a deviation from upper triangular matrix as shown below
\begin{align}
    \begin{bmatrix}
        a_0 & & & &  \\
         & a_1 & & &  \\
         & & a_2 & & \\
         & &  & a_3 &a_4\\
         & & & a_5& a_6\\
    \end{bmatrix}
\end{align}
To "fix" this problem, we can stop the iterations at some place and check for these blocks and check for these blocks and solve for the eigenvalues of this trivial $2 \times 2$ manually.\\

An elegant solution to the problem comes through this notion of \textbf{shifts} and \textbf{deflation} which will be explored in the next chapter
 
% Chapter Template

\chapter{Rayleigh Quotients and Wilkinson Shifts}

\label{Chapter5} % Change X to a consecutive number

\lhead{Chapter 5. \emph{Rayleigh Quotients and Wilkinson Shift}}

\section{Why Shift?}
As discussed in previous chapter, the convergence of Hessenberg QR algorithm is proportional to
$\frac{\lambda_i}{\lambda_j}$. Suppose if the eigenvalues are really close, in that case the algorithm will converge arbitrarily slowly. A way to get around this problem is to take a number
$\mu$ and perform the QR iterations on $H - \mu I $and add back $\mu I$ to the new RQ matrix. Precisely
\begin{algorithm}
    \caption{Demonstration of Shifts}
    \begin{algorithmic}
        \STATE $H_k - \mu I$:= QR
        \STATE $H_{k+1} = RQ + \mu I$
    \end{algorithmic}
\end{algorithm}

What this accomplishes is the following:\\ The rate of convergence is now proportional to 
\begin{align}
    \frac{|\lambda_i - \mu|}{|\lambda_j - \mu|}\\
\end{align}
This increases the rate of convergence and the amount by which the rate increases depends on how close we can get to $\lambda_j$. The problem is that we don't know the eigenvalues (obviously). So the agenda now is to find a way to come up with an on-the-fly way to approximate the eigenvalues in between the iterations

\section{Rayleigh Quotients}
\subsection{The Method}
One such heuristic is the Rayleigh quotient shift: Set the shift $\sigma_k$ in the $k^{th}$ step of the QR algorithm equal to the last diagonal element:
\begin{align}
    \sigma_k := h_{n,n}^{(k-1})
\end{align}
\begin{algorithm}
\caption{The Hessenberg QR Algorithm with Rayleigh Quotient Shift}
\begin{algorithmic}[1]
\STATE Let \(H_0 = H \in \mathbb{C}^{n \times n}\) be an upper Hessenberg matrix. This algorithm computes its Schur normal form \(H = UTU^*\).
\STATE Initialize \(k := 0\).
\FOR{\(m = n, n-1, \dots, 2\)}
    \REPEAT
        \STATE \(k := k + 1\).
        \STATE Compute the shift: \(\sigma_k := h^{(k-1)}_{m,m}\).
        \STATE Factorize \(H_{k-1} - \sigma_k I = Q_kR_k\).
        \STATE Update \(H_k := R_kQ_k + \sigma_k I\).
        \STATE Update \(U_k := U_{k-1}Q_k\).
    \UNTIL{\(|h^{(k)}_{m,m-1}|\) is sufficiently small.}
\ENDFOR
\STATE Set \(T := H_k\).
\end{algorithmic}
\end{algorithm}

Algorithm 4.4 implements this heuristic. Notice that the shift changes in each iteration step! Notice also that deflation is incorporated in Algorithm 4.4. As soon as the last lower off-diagonal element is sufficiently small, it is declared zero, and the algorithm proceeds with a smaller matrix. In Algorithm 4.4 the "active portion" of the matrix is $m \times m$. One thing which I won't elaborate much on is that implementing Rayleigh shifts achieves what is called quadratic convergence, as compared to the linear convergence of the naive QR algorithm. 
\subsection{Where Does it Fail?}
We still haven't addressed the issue of equal eigenvalues. To illustrate this, consider the following matrix
\begin{align}
    A = \begin{bmatrix}
        0 & 1\\
        1 & 0\\
    \end{bmatrix}
\end{align}
Here the value of $\sigma_k$ will be stuck at 0 since the eigenvalues are -1 and 1 and the algorithm can't choose which way to go, so it slowly decays to nothingness, leaving us with an infinite loop. This is the true reason why Jordan blocks show up. This problem showed up due a symmetry between the eigenvalues and $\sigma_k$. The next method breaks this symmetry and gives us a final method which will converge regardless of the entries in the matrix.

\section{Wilkinson Shift}
The heuristic chosen by this version is,
\begin{align}
    \sigma_k := \text{eigenvalue of} \begin{bmatrix}
        h_{n-1, n-1}^{(k-1)} & h_{n-1, n}^{(k-1)}\\
        h_{n, n-1}^{(k-1)} & h_{n, n}^{(k-1)}\\
    \end{bmatrix}
\end{align}
Implementing this gives us our ultimate goal. An algorithm which produces eigenvalues of any arbitrary matrix, complex or not. (Also, as an added bonus it also achieves cubic convergence)

% Chapter Template

\chapter{Final Form of the Algorithm} % Main chapter title

\label{Chapter6}

\lhead{Chapter 6. \emph{Final Form of the Algorithm}} % Change X to a consecutive number; this is for the header on each page - perhaps a shortened title

%----------------------------------------------------------------------------------------
%	SECTION 1
%----------------------------------------------------------------------------------------

\begin{algorithm}
\caption{Eigenvalue Calculation}
\begin{algorithmic}
\STATE To Upper Hessenberg form (via Householder reflection):
    \FOR{$k = 1$ TO $n-2$}
        \STATE Compute Householder reflector $P_k$ 
        \STATE Update $A$: $A = P_k A P_k^T$
    \ENDFOR

\STATE QR Iterations (via Givens Rotation):
\WHILE{$h_{n,n} > \epsilon$}
        \FOR{$k = 1$ TO $n-1$}
            \STATE Compute Givens rotation matrix $G_k$ to zero out the $(k+1, k)$ element.
            \STATE Apply $G_k$ from the left and right to $A$.
        \ENDFOR
    \STATE Form $A = R Q$
\ENDWHILE
\STATE Extract Eigenvalues from Quasi-Triangular Matrix
\FOR{$i = 1$ TO $n$}
        \IF{$|A_{i+1, i}| < \epsilon$}
            \STATE Eigenvalue: $\lambda_i = A_{i, i}$
        \ELSE
            \STATE Compute eigenvalues of the $2 \times 2$ sub-matrix by solving quadratic:
            \begin{align*}
            \begin{bmatrix}
                A_{i, i} & A_{i, i+1} \\
                A_{i+1, i} & A_{i+1, i+1}
            \end{bmatrix}
            \end{align*}
        \ENDIF
    \ENDFOR
\RETURN Eigenvalues $\lambda_1, \lambda_2,\cdots, \lambda_n$
\end{algorithmic}
\end{algorithm}








% Chapter Template

\chapter{Comparison With Other Algorithms} % Main chapter title

\label{Chapter6}

\lhead{Chapter 6. \emph{Comparison With Other Algorithms}} % Change X to a consecutive number; this is for the header on each page - perhaps a shortened title

%----------------------------------------------------------------------------------------
%	SECTION 1
%----------------------------------------------------------------------------------------

In true Indian fashion, we shall compare our little algorithm with its cousins and debate which one is the best. Below is a table\footnote{Credits: Wikipedia} which outlines pros and cons of other eigenalgorithms on the market.

\begin{longtable}{|p{2.5cm}|p{3.5cm}|p{3.5cm}|p{2.5cm}|p{2.5cm}|}
\hline
\textbf{Algorithm} & \textbf{Pros} & \textbf{Cons} & \textbf{Time Complexity} & \textbf{Conditions} \\ \hline
\endfirsthead
\hline
\textbf{Algorithm} & \textbf{Pros} & \textbf{Cons} & \textbf{Time Complexity} & \textbf{Conditions} \\ \hline
\endhead
Householder + Givens & Stable, precise; preserves sparsity. & Complex preprocessing for Hessenberg form. & \( O(n^3) \) & Any square matrix. \\ \hline
Jacobi Method & Accurate for symmetric matrices; simple. & Slow for non-symmetric matrices; inefficient for large ones. & \( O(n^3) \) (dense) & Symmetric (real/complex). \\ \hline
Lanczos Algorithm & Efficient for sparse symmetric matrices. & Sensitive to rounding; limited to Hermitian matrices. & \( O(kn^2) \) & Sparse, symmetric/Hermitian. \\ \hline
QR Algorithm & Robust for eigenvalue computation. & High cost for large dense matrices. & \( O(n^3) \) & Any square matrix. \\ \hline
Arnoldi Algorithm & Handles non-symmetric sparse matrices. & Requires re-orthogonalization; less stable. & \( O(kn^2) \) & Sparse, any matrix. \\ \hline
Power Iteration & Simple, finds dominant eigenvalue. & Only computes the largest eigenvalue; slow for close eigenvalues. & \( O(kn^2) \) & Any square matrix. \\ \hline
\end{longtable}

In conclusion, the algorithm which we chose was the only true algorithm which can compute all eigenvalues of an arbitrary matrix in a time which is less than expected lifetime of the universe.


% Chapter Template

\chapter{Bibliography} % Main chapter title

\label{Bibliography}

\lhead{Bibliography. \emph{Bibliography}} % Change X to a consecutive number; this is for the header on each page - perhaps a shortened title

%----------------------------------------------------------------------------------------
%	SECTION 1
%----------------------------------------------------------------------------------------
\begin{enumerate}
    \item \href{https://people.inf.ethz.ch/arbenz/ewp/Lnotes/chapter4.pdf}{A Handout on Eigenvalue Algorithms, by ETH Zurich}
    \item \href{https://www.math.iit.edu/~fass/477577_Chapter_11.pdf}{A Handout on Convergence of QR Algorithm, Illinois Institute of Technology}
    \item \href{https://dspace.mit.edu/bitstream/handle/1721.1/75282/18-335j-fall-2006/contents/lecture-notes/lec16.pdf}{Lecture 16 on Shifted QR Algorithms, Massachusets Institute of Technology}
    \item \href{https://dl.acm.org/doi/pdf/10.1145/800256.810675#:~:text=The%20QR%20Algorithm%20converges%20for%20an%20unreduced%20Hessenberg%20matrix%20if,and%20two%20o%2Fodd%20multiplicity.}{Necessary and Sufficient Conditions for Convergence of QR Algorithm on Hessenberg Matrix, BERESFORD PARLETT, University of California Berkeley, California}
\end{enumerate}

%----------------------------------------------------------------------------------------
%	THESIS CONTENT - APPENDICES
%----------------------------------------------------------------------------------------

\addtocontents{toc}{\vspace{2em}} % Add a gap in the Contents, for aesthetics

%\appendix % Cue to tell LaTeX that the following 'chapters' are Appendices

% Include the appendices of the thesis as separate files from the Appendices folder
% Uncomment the lines as you write the Appendices

%% Appendix Template

\chapter{Local plastic strain measurements} % Main appendix title

\label{AppendixX} % Change X to a consecutive letter; for referencing this appendix elsewhere, use \ref{AppendixX}

\lhead{Appendix A. \emph{Local plastic strain measurements}} % Change X to a consecutive letter; this is for the header on each page - perhaps a shortened title


The plastic deformation in metallic materials is of a heterogeneous nature, with strain distributed non-uniformly at the inter and intragranular levels. In order to evaluate the capabilities of full-field models for predicting the deformation arrangement, experimental techniques are required to measure the spatially resolved deformation formation in microstructures.

\section{Electron Backscatter Diffraction}
\label{Electron Backscatter Diffraction}

Electron Backscatter Diffraction (EBSD) is a technique for measuring crystal orientation using a scanning electron microscope (SEM). A focussed electron beam illuminates a spot on the surface of the sample that is inclined by $70^{\circ}$  to the beam direction. Part of this primary beam is elastically scattered from the crystal lattice. These backscattered electrons diffract through the crystal lattice as they escape the material, producing a spatial electron intensity pattern that indicates the crystal structure and orientation. The diffracted electors emanate spherically from the illuminated point on the sample surface. Since the sample is tilted, a detector can be positioned to collect a portion of the diffraction pattern on a flat plane. The collected Kikuchi pattern (named for a pioneer in electron diffraction) can then be analysed to assess the orientation of the crystal by measuring the position and orientation of bands of constructive interference \cite{maitland2007electron}.

\noindent A crystal orientation map can be collected by rastering the beam across the sample surface, collecting and analysing a Kikuchi pattern at each point. These maps contain information about grain size, grain boundary character, sample texture and alloy phase. EBSD can also be used to assess local deformations \cite{wright2011review}. This is achieved either by direct measurement of elastic strains through cross-correlation of individual Kikuchi patterns \cite{wilkinson2006high,wilkinson2006high,britton2012high} or by quantifying small changes and orientation gradients that are generated by plastic deformation \cite{kamaya2004measurement,kamaya2007local,githinji2013ebsd}.

\subsection{Measures of misorientation}
\label{Measures of misorientation}

Misorientation is the difference between two orientations and can be expressed as the minimum angle of rotation around any axis that transforms one orientation into another. The plastic deformation of a crystal due to dislocation slip induces a rotation of the crystal lattice. In an unrestrained single crystal, this leads to a net rotation of the material. However, in order to ensure that the grains remain compatible with their neighbours in a polycrystal, deformation gradients and differences in active slip system within single grains are required \cite{barbe2001intergranular}. This leads to orientation gradients that develop within the grains. Each active slip system induces rotation around a different axis, and the amount of rotation is proportional to the amount of slip. The changes in orientation are accommodated in the crystal lattice by means of geometrically necessary dislocations (GNDs), which cause a net rotation of the lattice. A GND density can be calculated from the rotation gradients of the lattice \cite{pantleon2008resolving}.

\noindent The measurement parameters for misorientation define between which orientations the misorientation is calculated and how this relates to the crystal rotation. In a class of misorientation parameters, a misorientation between neighbouring spatial points of an EBSD is calculated. This could be a single point in a single direction or, more generally, in a kernel of positions surrounding each point. The average of the misorientation at each point in the kernel results in the kernel-averaged misorientation (KAM) \cite{wright2011review}. The KAM quantifies local fluctuations in orientation. Large values indicate a large local orientation gradient resulting from a change in the active slip system or an abrupt change in slip activity. Examples of KAM maps taken from a plastically deformed sample are shown in figure \ref{fig:Measures of misorientation1}a-b. The EBSD maps are taken at two different EBSD step sizes and show that KAM values are sensitive to the chosen EBSD step size.

\begin{figure}
    \centering
    \includegraphics[scale=1.5]{Pictures/KAM.jpg}
    \caption{Examples of misorientation parameters for the same region of material. a) and b) are KAM calculated for a 3x3 pixel $^{2}$ kernel for EBSD step sizes of $0.2 \mu \mathrm{m}$ and $0.4 \mu \mathrm{m} .$ c) is ROD calculated to each grains average orientation and d) ROD calculated to the point in each grain with minimum KAM. Taken from \cite{wright2011review}.}
    \label{fig:Measures of misorientation1}
\end{figure}

KAM does not indicate gradual changes of orientation relating to gradients of deformation. This can be measured using a second class of parameters, the reference orientation deviation (ROD) \cite{wright2011review}. As the name suggests ROD is the misorientation between each point of a map to a reference orientation.A reference orientation must be defined separately for each grain in the EBSD map, since the misorientation between two grains is generally greater than any gradient in it. Common choices for the reference orientation are the mean orientation of each grain or the point in each grain with the minimum KAM value. Figure \ref{fig:Measures of misorientation1}c-d shows a ROD map for each of these reference orientation choices. Similar gradients can be seen in each map but with different magnitudes of the ROD values, whereby the choice of the reference orientation only changes the zero point of the misorientation. Since ROD measures misorientation relative to a fixed reference, it is not sensitive to step size changes.

\noindent The use of misorientation can indicate intense strain in the grains of a microstructure \cite{kamaya2006quantification}, but cannot identify magnitudes of strains. The distribution of the intragranular plastic strain is also not well predicted due to misorientation compared to local strain measurements from image correlation techniques \cite{kamaya2007local}. The misorientation, however, provides a good measure of the microstructure evolution induced by plastic deformation.
%\input{Appendices/AppendixB}
%\input{Appendices/AppendixC}

\addtocontents{toc}{\vspace{2em}} % Add a gap in the Contents, for aesthetics

\backmatter

%----------------------------------------------------------------------------------------
%	BIBLIOGRAPHY
%----------------------------------------------------------------------------------------
 % The references (bibliography) information are stored in the file named "Bibliography.bib"

\end{document}  
